% 7-02(7쪽) — 보기 ①~⑤ 의 그래프 다섯 개(x=0 에서 극한값이 존재하는 그래프 고르기). 보기 자체가 그림이라 본문 아래에 둔다.
% 스캔 화질이 낮아 TikZ 로 다시 그림 — 라벨은 원문 그대로(O, x, y, y=f(x) · ⑤ 의 눈금 -2, -1, 1, 2, 3 과 1)만 두고 원문에 없는 눈금·값은 넣지 않는다.
\begin{tikzpicture}[x=0.5cm, y=0.5cm, line width=0.5pt, font=\footnotesize]
  % ① x→0- 에서 위로, x→0+ 에서 아래로 갈라지는 쌍곡선 꼴
  \begin{scope}
    \node at (-3.5,2.6) {①};
    \draw[->, >=stealth, line width=0.7pt] (-2.85,0) -- (2.9,0) node[below=1pt] {$x$};
    \draw[->, >=stealth, line width=0.7pt] (0,-2.9) -- (0,2.95) node[right=1pt] {$y$};
    \node[below left=-2pt and -3pt] at (0,0) {$\pt{O}$};
    \draw[line width=0.6pt, domain=-2.75:-0.38, samples=90, smooth] plot (\x, {-1/\x});
    \draw[line width=0.6pt, domain=0.38:2.75, samples=90, smooth] plot (\x, {-1/\x});
    \node[anchor=east] at (-0.45,2.6) {$y=f(x)$};
  \end{scope}
  % ② 양쪽 모두 아래로 발산하는 곡선
  \begin{scope}[shift={(8.4,1.9)}]
    \node at (-3.5,0.6) {②};
    \draw[->, >=stealth, line width=0.7pt] (-2.85,0) -- (2.9,0) node[above=1pt] {$x$};
    \draw[->, >=stealth, line width=0.7pt] (0,-3.2) -- (0,1.05) node[left=1pt] {$y$};
    \node[below left=-2pt and -3pt] at (0,0) {$\pt{O}$};
    \draw[line width=0.6pt, domain=-2.75:-0.52, samples=90, smooth] plot (\x, {-0.8/(\x*\x)});
    \draw[line width=0.6pt, domain=0.52:2.75, samples=90, smooth] plot (\x, {-0.8/(\x*\x)});
    \node[anchor=west] at (0.7,-2.75) {$y=f(x)$};
  \end{scope}
  % ③ x=0 에 빈 점이 있고 그 아래 채운 점이 있는 증가 곡선
  \begin{scope}[shift={(0,-7.4)}]
    \node at (-3.5,2.6) {③};
    \draw[->, >=stealth, line width=0.7pt] (-2.85,0) -- (2.9,0) node[below=1pt] {$x$};
    \draw[->, >=stealth, line width=0.7pt] (0,-2.9) -- (0,2.95) node[left=1pt] {$y$};
    \node[below left=-2pt and -3pt] at (0,0) {$\pt{O}$};
    \draw[line width=0.6pt, domain=-1.5:1.35, samples=90, smooth] plot (\x, {0.9*\x*\x*\x+0.8});
    \draw[fill=white, line width=0.5pt] (0,0.8) circle (2pt);
    \fill (0,-0.8) circle (2pt);
    \node[anchor=west] at (1.3,2.7) {$y=f(x)$};
  \end{scope}
  % ④ x=0 에서 좌우가 갈라지는 두 직선
  \begin{scope}[shift={(8.4,-7.4)}]
    \node at (-3.5,2.6) {④};
    \draw[->, >=stealth, line width=0.7pt] (-2.85,0) -- (2.9,0) node[below=1pt] {$x$};
    \draw[->, >=stealth, line width=0.7pt] (0,-2.9) -- (0,2.95) node[left=1pt] {$y$};
    \node[below left=-2pt and -3pt] at (0,0) {$\pt{O}$};
    \draw[line width=0.6pt] (-1.3,3.02) -- (0,1.3);
    \draw[line width=0.6pt] (0,-0.6) -- (1.4,-2.45);
    \draw[fill=white, line width=0.5pt] (0,1.3) circle (2pt);
    \draw[fill=white, line width=0.5pt] (0,-0.6) circle (2pt);
    \node[anchor=west] at (0.25,1.3) {$y=f(x)$};
  \end{scope}
  % ⑤ 한 칸마다 올라가는 톱니 꼴(각 구간 왼쪽 끝은 채운 점 · 오른쪽 끝은 빈 점)
  \begin{scope}[shift={(0,-14.2)}]
    \node at (-3.5,1.9) {⑤};
    \draw[->, >=stealth, line width=0.7pt] (-3.1,0) -- (4.1,0) node[below=1pt] {$x$};
    \draw[->, >=stealth, line width=0.7pt] (0,-1.8) -- (0,2.4) node[left=1pt] {$y$};
    \draw[densely dotted, line width=0.4pt] (-1,1) -- (3,1);
    \foreach \n in {-1,0,1,2,3} { \draw[densely dotted, line width=0.4pt] (\n,1) -- (\n,0); }
    \foreach \n in {-2,-1,0,1,2} { \draw[line width=0.6pt] (\n,0) -- (\n+1,1); }
    \foreach \n in {-2,-1,0,1,2} { \fill (\n,0) circle (2pt); }
    \foreach \n in {-1,0,1,2,3} { \draw[fill=white, line width=0.5pt] (\n,1) circle (2pt); }
    \node[above left=-2pt and -1pt] at (0,1) {$1$};
    \node[below=1pt] at (-2,0) {$-2$};
    \node[below=1pt] at (-1,0) {$-1$};
    \node[below left=-2pt and -3pt] at (0,0) {$\pt{O}$};
    \node[below=1pt] at (1,0) {$1$};
    \node[below=1pt] at (2,0) {$2$};
    \node[below=1pt] at (3,0) {$3$};
    \node[anchor=west] at (0.8,1.95) {$y=f(x)$};
  \end{scope}
\end{tikzpicture}
